\section{洛伦兹变换}
\subsection{一般形式}
在闵可夫斯基空间中，度规张量为：
\begin{equation*}
	\vb{\eta} = \eta_{\mu\nu} = \eta^{\mu\nu} = \begin{pmatrix}
		-1 & 0 & 0 & 0 \\
		0  & 1 & 0 & 0 \\
		0  & 0 & 1 & 0 \\
		0  & 0 & 0 & 1
	\end{pmatrix}.
\end{equation*}

若坐标变换 \(x^\mu \to \tilde{x}^\mu = \Lambda^\mu{}_\nu x^\nu\) 保持时空间隔不变，即
\begin{equation*}
	\d s^2 = -\eta_{\mu\nu} x^\mu x^\nu = -\tilde{\eta}_{pq} \tilde{x}^p \tilde{x}^q = -\eta_{pq} \Lambda^\mu{}_p \Lambda^\nu{}_q x^\mu x^\nu.
\end{equation*}
则要求 \(\eta_{pq} \Lambda^\mu{}_p \Lambda^\nu{}_q = \eta_{\mu\nu}\).

先考虑无限小变换情形 \(\Lambda^p{}_q = \delta^p{}_q + \chi^p{}_q\)，\(|\chi^p{}_q| \ll 1\)，代入得，也即
\begin{equation*}
	\eta_{p\nu} \chi^\mu{}_p = -\eta_{\mu q} \chi^\nu{}_q \quad\implies\quad\vb{\eta} \vb{\chi} \vb{\eta} = -\vb{\chi}^\mathrm{T}.
\end{equation*}
解得，\(\chi\) 有6个自由度（用\(\vb{\varphi},\,\vb{\theta}\)表征），其形式为
\begin{align*}
	 & \vb{\chi} = -\vb{\varphi} \cdot \vb{K} - \vb{\theta} \cdot \vb{S}                                               \\
	 & K_1 = \begin{pmatrix}
		         0 & 1 & 0 & 0 \\
		         1 & 0 & 0 & 0 \\
		         0 & 0 & 0 & 0 \\
		         0 & 0 & 0 & 0
	         \end{pmatrix}, \quad
	K_2 = \begin{pmatrix}
		      0 & 0 & 1 & 0 \\
		      0 & 0 & 0 & 0 \\
		      1 & 0 & 0 & 0 \\
		      0 & 0 & 0 & 0
	      \end{pmatrix}, \quad
	K_3 = \begin{pmatrix}
		      0 & 0 & 0 & 1 \\
		      0 & 0 & 0 & 0 \\
		      0 & 0 & 0 & 0 \\
		      1 & 0 & 0 & 0
	      \end{pmatrix}.                                                                                               \\
	 & S_1 = \begin{pmatrix}
		         0 & 0 & 0 & 0  \\
		         0 & 0 & 0 & 0  \\
		         0 & 0 & 0 & -1 \\
		         0 & 0 & 1 & 0
	         \end{pmatrix}, \quad
	S_2 = \begin{pmatrix}
		      0 & 0  & 0 & 0 \\
		      0 & 0  & 0 & 1 \\
		      0 & 0  & 0 & 0 \\
		      0 & -1 & 0 & 0
	      \end{pmatrix}, \quad
	S_3 = \begin{pmatrix}
		      0 & 0 & 0  & 0 \\
		      0 & 0 & -1 & 0 \\
		      0 & 1 & 0  & 0 \\
		      0 & 0 & 0  & 0
	      \end{pmatrix}.                                                                                               \\
	 & [S_i, S_j] = \epsilon_{ijk} S_k, \quad [S_i, K_j] = \epsilon_{ijk} K_k, \quad [K_i, K_j] = -\epsilon_{ijk} S_k.
\end{align*}

进而，易知
\begin{equation*}
	\vb{\chi} = -\vb{\varphi} \cdot \vb{K} - \vb{\theta} \cdot \vb{S} = \begin{pmatrix}
		0             & -\vb{\varphi}^\mathrm{T}        \\
		-\vb{\varphi} & \vb{\theta} \cdot \vb{\epsilon}
	\end{pmatrix}, \quad
	[\vb{\varphi} \cdot \vb{K}, \vb{\theta} \cdot \vb{S}] = (\vb{\theta} \times \vb{\varphi}) \cdot \vb{K}.
\end{equation*}
其中 \(\vb{\epsilon}\) 是三阶全反对称张量（\(\epsilon_{123} = 1\)，\(\epsilon_{132} = -1\)，…），点乘对最临近指标缩并.
\begin{notes}
	本节所有矢量均视为3×1的矩阵，并使用转置符号\(\mathrm{T}\)以方便从矩阵视角进行运算，当然，为了清晰起见，在一些地方仍用 \(\vb{f} \cdot \vb{g}\) 而非 \(\vb{f}^\mathrm{T} \vb{g}\) 来表示矢量的点积（一些张量的缩并也同样用更直观的点积形式）.
\end{notes}

正常量级的变换可看为\(N\)次（\(N \to \infty\)）连续的无限小变换，将无穷小变换中的 \(\vb{\varphi}\)，\(\vb{\theta}\) 改写为 \(\vb{\varphi}/N\)，\(\vb{\theta}/N\)（\(\vb{\varphi}\)，\(\vb{\theta}\) 具有正常量级），得：
\begin{equation*}
	\Lambda(\vb{\varphi}, \vb{\theta}) = \left( \bm{\mathcal{I}} - \dfrac{\vb{\varphi}}{N} \cdot \vb{K} - \dfrac{\vb{\theta}}{N} \cdot \vb{S} \right)^N = \exp\left( -\vb{\varphi} \cdot \vb{K} - \vb{\theta} \cdot \vb{S} \right).
\end{equation*}
其中 \(\bm{\mathcal{I}}\) 是4×4单位矩阵（注意与3×3单位矩阵\(\vb{I}\)区分）.

以上是保持时空间隔不变下，变换矩阵应具有的形式，其中包括纯空间转动项 \(\exp(-\vb{\theta} \cdot \vb{S})\)（对应绕 \(\uv{\theta}\) 方向的轴旋转\(\theta\)弧度），以及Boost项 \(B(\vb{\varphi}) = \exp(-\vb{\varphi} \cdot \vb{K})\).

对于 \(n \in \mathbb{Z}^+\)，易知：
\begin{equation*}
	(\vb{\varphi} \cdot \vb{K})^{2n} = \begin{pmatrix}
		0                      & \varphi \uv{e}_\varphi^\mathrm{T} \\
		\varphi \uv{e}_\varphi & \vb{0}
	\end{pmatrix}^{2n} = \varphi^{2n} \begin{pmatrix}
		1      & \vb{0}                                   \\
		\vb{0} & \uv{e}_\varphi \uv{e}_\varphi^\mathrm{T}
	\end{pmatrix},\quad
	(\vb{\varphi} \cdot \vb{K})^{2n-1} = \varphi^{2n-1} \begin{pmatrix}
		0              & \uv{e}_\varphi^\mathrm{T} \\
		\uv{e}_\varphi & \vb{0}
	\end{pmatrix}.
\end{equation*}
进而
\begin{equation*}
	B(\vb{\varphi}) = \exp(-\vb{\varphi} \cdot \vb{K}) = \sum_{k=0}^\infty \dfrac{(\vb{\varphi} \cdot \vb{K})^k}{k!} = \begin{pmatrix}
		\cosh\varphi                    & -\sinh\varphi \, \uv{e}_\varphi^\mathrm{T}                           \\
		-\sinh\varphi \, \uv{e}_\varphi & \vb{I} + (\cosh\varphi - 1) \uv{e}_\varphi \uv{e}_\varphi^\mathrm{T}
	\end{pmatrix}.
\end{equation*}
记 \(\vb{\varphi} = \bm{\hat{e}_\beta} \artanh\beta,\,\gamma = (1 - \beta^2)^{-1/2}\)，其中\(0 \leq \beta < 1\)，则：
\begin{equation*}
	M(\vb{\beta}) = B(\vb{\varphi}) = \exp(-\vb{\varphi} \cdot \vb{K}) = \exp\left( -\vb{\beta} \cdot \vb{K}\, \dfrac{\artanh\beta}{\beta} \right) = \begin{pmatrix}
		\gamma                & -\gamma \, \vb{\beta}^\mathrm{T}                                       \\
		-\gamma \, \vb{\beta} & \vb{I} + \dfrac{\gamma^2}{\gamma + 1} \vb{\beta} \vb{\beta}^\mathrm{T}
	\end{pmatrix}.
\end{equation*}

\subsection{变换实例}
以下不考虑整体的转动，只考虑矢量、张量在Boost矩阵影响下的变换.
\subsubsection{矢量的变换}

对于矢量，其变换为（以下 \(\vb{\beta}\) 系中的物理量带右上标\('\)，地面系则无右上标）：
\begin{equation*}
	\begin{pmatrix}
		(A^0)' \\
		\vb{A}'
	\end{pmatrix}
	=
	\begin{pmatrix}
		\gamma                & -\gamma \, \vb{\beta}^\mathrm{T}                                       \\
		-\gamma \, \vb{\beta} & \vb{I} + \dfrac{\gamma^2}{\gamma + 1} \vb{\beta} \vb{\beta}^\mathrm{T}
	\end{pmatrix}
	\begin{pmatrix}
		A^0 \\
		\vb{A}
	\end{pmatrix}.
\end{equation*}

展开得：
\begin{equation*}
	(A^0)' = \gamma A^0 - \gamma \, \vb{\beta} \cdot \vb{A}, \quad
	\vb{A}' = -\gamma \, \vb{\beta} \, A^0 + \vb{A} + \dfrac{\gamma^2}{\gamma + 1} \vb{\beta} \left( \vb{\beta} \cdot \vb{A} \right).
\end{equation*}

从张量理论或直接计算的角度均可说明，四维矢量的缩并为不变量：
\begin{equation*}
	A^\mu A_\mu = (A^0)^2 - |\vb{A}|^2 = (A^0\,')^2 - |\vb{A}'|^2.
\end{equation*}

\begin{example}
	四维坐标 \(x^\mu = (\vb{r}, \, ct)^\mathrm{T} \to (\vb{r}', \, ct')^\mathrm{T}\)：
	\begin{equation*}
		ct' = \gamma \, ct - \gamma \, \vb{\beta} \cdot \vb{r}, \quad
		\vb{r}' = -\gamma \, \vb{\beta} \, ct + \vb{r} + \dfrac{\gamma^2}{\gamma + 1} \vb{\beta} \left( \vb{\beta} \cdot \vb{r} \right).
	\end{equation*}
	四维坐标缩并 \(x^\mu x_\mu = -c^2 t^2 + r^2\) 为不变量.
\end{example}

\begin{example}
	四维速度 \(U^\mu = \d x^\mu/\d \tau = \left( \gamma_u c, \, \gamma_u \vb{\beta}_u c \right)^\mathrm{T} \to \left( \gamma_{u'} c, \, \gamma_{u'} \vb{\beta}_{u'} c \right)^\mathrm{T}\)：
	\begin{equation*}
		\vb{\beta}_{u'} = \dfrac{ \vb{\beta}_u - \gamma \, \vb{\beta} + \dfrac{\gamma^2}{\gamma + 1} \vb{\beta} \left( \vb{\beta} \cdot \vb{\beta}_u \right) }{ \gamma \left( 1 - \vb{\beta} \cdot \vb{\beta}_u \right) }, \quad
		\gamma_{u'} = \gamma \gamma_u \left( 1 - \vb{\beta} \cdot \vb{\beta}_u \right).
	\end{equation*}
	四维速度缩并 \(U^\mu U_\mu = -c^2\) 为不变量.
\end{example}

\begin{example}
	对于四维加速度 \(A^\mu = \d u^\mu/\d \tau = \left( \gamma_u^4 \, \vb{a} \cdot \vb{\beta}_u, \, \gamma_u^4 \left( \vb{a} + \vb{\beta}_u \left( \vb{\beta}_u \cdot \vb{a} \right) \right) \right)\)：
	\begin{equation*}
		\vb{a}' = \dfrac{ \vb{a} \left( 1 - \vb{\beta} \cdot \vb{\beta}_u \right) + \left( \vb{a} \cdot \vb{\beta} \right) \left( \vb{\beta}_u - \dfrac{\gamma}{\gamma + 1} \vb{\beta} \right) }{ \gamma^2 \left( 1 - \vb{\beta} \cdot \vb{\beta}_u \right) }.
	\end{equation*}
	四维加速度缩并 \(A^\mu A_\mu = \gamma_u^6 \left( a^2 - \left( \vb{a} \cdot \vb{\beta}_u \right)^2 \right)\) 为不变量.
\end{example}

\begin{example}
	对于四维电流密度 \(J^\mu = (\rho c, \, \vb{J})\)，设 \(\vb{\beta}\) 系恰为导体系，有 \(\vb{J}' = \sigma \vb{E}'\)，则：
	\begin{equation*}
		\vb{J} = \rho \, \vb{\beta} \, c + \gamma \sigma \left( \vb{E} + \vb{\beta} \times \vb{B} \, c - \vb{\beta} \left( \vb{\beta} \cdot \vb{E} \right) \right), \quad
		\rho = \gamma \left( \rho' + \dfrac{\sigma}{c} \, \vb{\beta} \cdot \vb{E} \right).
	\end{equation*}
	若导体内无自由电荷，即 \(\rho' = 0\)，则：
	\begin{equation*}
		\vb{J} = \gamma \sigma \left( \vb{E} + \vb{\beta} \times \vb{B} \, c \right), \quad
		\rho = \dfrac{\gamma \sigma}{c} \, \vb{\beta} \cdot \vb{E}.
	\end{equation*}
\end{example}

\subsubsection{二阶张量的变换}

对于二阶混合张量，其变换为：
\begin{equation*}
	\begin{pmatrix}
		w'      & \vb{S}^\mathrm{'T} \\
		\vb{G}' & \vb{\Sigma}'
	\end{pmatrix}
	=
	\begin{pmatrix}
		\gamma                & -\gamma \, \vb{\beta}^\mathrm{T}                                       \\
		-\gamma \, \vb{\beta} & \vb{I} + \dfrac{\gamma^2}{\gamma + 1} \vb{\beta} \vb{\beta}^\mathrm{T}
	\end{pmatrix}
	\begin{pmatrix}
		w      & \vb{S}^\mathrm{T} \\
		\vb{G} & \vb{\Sigma}
	\end{pmatrix}
	\begin{pmatrix}
		\gamma                & -\gamma \, \vb{\beta}^\mathrm{T}                                       \\
		-\gamma \, \vb{\beta} & \vb{I} + \dfrac{\gamma^2}{\gamma + 1} \vb{\beta} \vb{\beta}^\mathrm{T}
	\end{pmatrix}^{-1}.
\end{equation*}

展开得：
\begin{align*}
	w' \;           & =\; \gamma^{2}\Big( w + \vb{S}\!\cdot\!\vb{\beta}
	-\vb{G}\!\cdot\!\vb{\beta}
	-\vb{\beta}\!\cdot\!\vb{\Sigma}\!\cdot\!\vb{\beta} \Big),                     \\[4pt]
	\vb{S}' \;      & =\; \gamma^{2} w\,\vb{\beta}
	+\gamma\,\vb{S}
	+\frac{\gamma^{3}}{\gamma+1}\,\vb{\beta}\,(\vb{\beta}\!\cdot\!\vb{S})
	-\gamma^{2}\,\vb{\beta}\,(\vb{\beta}\!\cdot\!\vb{G})
	-\gamma\,(\vb{\beta}\!\cdot\!\vb{\Sigma})
	-\frac{\gamma^{3}}{\gamma+1}\,\vb{\beta}\,(\vb{\beta}\vb{\beta}:\vb{\Sigma}), \\[4pt]
	\vb{G}' \;      & =\; -\gamma^{2} w\,\vb{\beta}
	+\gamma\,\vb{G}
	+\frac{\gamma^{3}}{\gamma+1}\,\vb{\beta}\,(\vb{\beta}\!\cdot\!\vb{G})
	-\gamma^{2}\,\vb{\beta}\,(\vb{\beta}\!\cdot\!\vb{S})
	+\gamma\,(\vb{\Sigma}\!\cdot\!\vb{\beta})
	+\frac{\gamma^{3}}{\gamma+1}\,\vb{\beta}\,(\vb{\beta}\vb{\beta}:\vb{\Sigma}), \\[4pt]
	\vb{\Sigma}' \; & =\; \vb{\Sigma}
	-\gamma^{2} w\,\vb{\beta}\vb{\beta}^{T}
	-\gamma\big(\vb{\beta}\vb{S}^{T}-\vb{G}\vb{\beta}^{T}\big)
	-\frac{\gamma^{3}}{\gamma+1}\,\vb{\beta}\vb{\beta}^{T}
	\big(\vb{\beta}\!\cdot\!\vb{S}-\vb{\beta}\!\cdot\!\vb{G}\big)                 \\
	                & \quad
	+\frac{\gamma^{2}}{\gamma+1}\Big(\vb{\Sigma}\!\cdot\!\vb{\beta}\vb{\beta}^{T}
	+\vb{\beta}\,(\vb{\beta}\!\cdot\!\vb{\Sigma})^{T}\Big)
	+\frac{\gamma^{4}}{(\gamma+1)^{2}}\,\vb{\beta}\vb{\beta}^{T}\,
	(\vb{\beta}\vb{\beta}:\vb{\Sigma}).
\end{align*}

\begin{notes}
	\(n \times n\) 方阵 \(M\) 的不变量

	记 \(M\) 的特征值为 \(\lambda_i\)（\(i = 1, 2, \dots, n\)），则：
	\begin{equation*}
		\det[\lambda I - M] = \sum_{k=0}^n e_k (-1)^k \lambda^{n-k}.
	\end{equation*}
	其中 \(e_k = \sum_{1 \leq i_1 < \dots < i_k \leq n} \lambda_{i_1} \cdots \lambda_{i_k}\)（\(k = 1, 2, \dots, n\)）构成 \(M\) 的 \(n\) 个不变量，\(e_0 = 1\).

	记
	\begin{equation*}
		p_k = \sum_{i=1}^n \lambda_i^k = \mathrm{tr}[M^k],\,(k \in \mathbb{Z}^+),\,E(t) = \sum_{k=0}^n e_k t^k = \prod_{i=1}^n (1 + \lambda_i t).
	\end{equation*}

	则
	\begin{equation*}
		\ln E(t) = \sum_{i=1}^n \ln(1 + \lambda_i t) = \sum_{i=1}^n \sum_{m=1}^\infty \dfrac{(-1)^{m-1}}{m} (\lambda_i t)^m = \sum_{m=1}^\infty \dfrac{(-1)^{m-1}}{m} p_m t^m.
	\end{equation*}

	进而
	\begin{equation*}
		\dfrac{\d E(t)}{\d t} = E(t) \dfrac{\d \ln E(t)}{\d t} = \sum_{j=0}^n e_j t^j \sum_{m=1}^\infty (-1)^{m-1} p_m t^{m-1} = \sum_{k=1}^\infty \sum_{j=0}^n p_{k-j} e_j t^k (-1)^{k-j-1}.
	\end{equation*}

	而 \(\d E(t)/\d t = \sum_{k=1}^n k e_k t^{k-1}\)，比较 \(t\) 的系数得：
	\begin{equation*}
		e_k = \frac{1}{k}\sum_{j=1}^n (-1)^{j-1} p_j e_{k-j}.
	\end{equation*}

	具体而言
	\begin{equation*}
		e_1 = \mathrm{tr}M,\,e_2 = \dfrac{1}{2}\left( (\mathrm{tr}M)^2 - \mathrm{tr}M^2 \right),\,e_3 = \dfrac{1}{6}\left( (\mathrm{tr}M)^3 - 3 \mathrm{tr}M \mathrm{tr}M^2 + 2 \mathrm{tr}M^3 \right),...,e_n = \det M.
	\end{equation*}

	另外可以证明，\(e_k\) 等于 \(M\) 的所有 \(k\) 阶主子式之和，但该性质不常用于具体计算.

\end{notes}

对于\(4\times4\)方阵，不变量为
\begin{equation*}
	e_1 = \mathrm{tr}M,\;e_2 = \dfrac{1}{2}\left( (\mathrm{tr}M)^2 - \mathrm{tr}M^2 \right),\;e_3 = \dfrac{1}{6}\left( (\mathrm{tr}M)^3 - 3 \mathrm{tr}M \mathrm{tr}M^2 + 2 \mathrm{tr}M^3 \right),\;e_4 = \det M.
\end{equation*}
矩阵写成如上的形式
\begin{equation*}
	M=\begin{pmatrix}
		w      & \vb{S}^\mathrm{T} \\
		\vb{G} & \vb{\Sigma}
	\end{pmatrix}.
\end{equation*}
则不变量进一步简化为
\begin{equation*}
	e_1 &= w + \mathrm{tr}\vb{\Sigma}, \quad e_2 = w \, \mathrm{tr}\vb{\Sigma} - \vb{S} \cdot \vb{G} + \dfrac{1}{2}\left( (\mathrm{tr}\vb{\Sigma})^2 - \mathrm{tr}\vb{\Sigma}^2 \right)\\
	e_3&=\det\vb{\Sigma}+\frac{w}{2}\Big((\operatorname{tr}\vb{\Sigma})^2-\operatorname{tr}(\vb{\Sigma}^2)\Big)
	+\mathbf{S}^{T}\vb{\Sigma}\mathbf{G}-(\operatorname{tr}\vb{\Sigma})(\mathbf{S}^{T}\mathbf{G})\\
	e_4&=w\det\vb{\Sigma}-\mathbf{S}^{T}\vb{\Sigma}^{2}\mathbf{G}
	+(\operatorname{tr}\vb{\Sigma})\,\mathbf{S}^{T}\vb{\Sigma}\mathbf{G}
	-\frac{1}{2}\Big((\operatorname{tr}\vb{\Sigma})^2-\operatorname{tr}(\vb{\Sigma}^2)\Big)(\mathbf{S}^{T}\mathbf{G})
\end{equation*}

\begin{example}
	对于 \((\vb{E}, \, \vb{B})\) 组合的电磁场张量，混合指标形式\(w = 0,\,\vb{\Sigma} = -\vb{\epsilon} \cdot \vb{B},\,\vb{S} = \vb{G} = \vb{E}/c\)（双上指标形式为\(w = 0,\,\vb{\Sigma} =- \vb{\epsilon} \cdot \vb{B},\,\vb{S} = -\vb{G} = \vb{E}/c\)），其不变量为
	\begin{equation*}
		e_1=e_3=0,\quad e_2=B^2-\frac{E^2}{c^2},\quad e_4=-\frac{\left(\vb{E}\cdot\vb{B}\right)^2}{c^2}
	\end{equation*}
	代入得变换后的混合指标形式为\(w' = 0,\,\vb{\Sigma}' =- \vb{\epsilon} \cdot \vb{B}',\,\vb{S}' = \vb{G}' = \vb{E}'/c\)，其中：
	\begin{equation*}
		\vb{E}' = \gamma \left( \vb{E} + \vb{\beta} \times c \, \vb{B} \right) - \dfrac{\gamma^2}{\gamma + 1} \vb{\beta} \left( \vb{\beta} \cdot \vb{E} \right),
	\end{equation*}
	\begin{equation*}
		\vb{B}' = \gamma \left( \vb{B} - \vb{\beta} \times \dfrac{\vb{E}}{c} \right) - \dfrac{\gamma^2}{\gamma + 1} \vb{\beta} \left( \vb{\beta} \cdot \vb{B} \right).
	\end{equation*}

	对于 \((\vb{D}, \, \vb{H})\) 或\((\vb{P}, \, \vb{M})\)组合的电磁场张量，同理有：
	\begin{equation*}
		\vb{D}' = \gamma \left( \vb{D} + \vb{\beta} \times \dfrac{\vb{H}}{c} \right) - \dfrac{\gamma^2}{\gamma + 1} \vb{\beta} \left( \vb{\beta} \cdot \vb{D} \right),
	\end{equation*}
	\begin{equation*}
		\vb{H}' = \gamma \left( \vb{H} - \vb{\beta} \times c \, \vb{D} \right) - \dfrac{\gamma^2}{\gamma + 1} \vb{\beta} \left( \vb{\beta} \cdot \vb{H} \right).
	\end{equation*}
	\begin{equation*}
		\vb{P}' = \gamma \left( \vb{P} + \vb{\beta} \times \dfrac{\vb{M}}{c} \right) - \dfrac{\gamma^2}{\gamma + 1} \vb{\beta} \left( \vb{\beta} \cdot \vb{P} \right),
	\end{equation*}
	\begin{equation*}
		\vb{M}' = \gamma \left( \vb{M} - \vb{\beta} \times c \, \vb{P} \right) - \dfrac{\gamma^2}{\gamma + 1} \vb{\beta} \left( \vb{\beta} \cdot \vb{M} \right).
	\end{equation*}

	若 \(\vb{\beta}\) 系恰为相对介电常数为 \(\varepsilon_r\)、相对磁导率为 \(\mu_r\) 的介质的本征系，即 \(\vb{B}' = \mu_0 \mu_r \vb{H}'\)，\(\vb{D}' = \varepsilon_0 \varepsilon_r \vb{E}'\)，则：
	\begin{align*}
		\vb{E} & = \dfrac{1}{\varepsilon_0 \varepsilon_r (1 - \beta^2)} \left[ (1 - \varepsilon_r \mu_r \beta^2) \vb{D} + (\varepsilon_r \mu_r - 1) \left( -\vb{\beta} \times \dfrac{\vb{H}}{c} + \vb{\beta} \left( \vb{\beta} \cdot \vb{D} \right) \right) \right], \\
		\vb{B} & = \dfrac{\mu_0}{\varepsilon_r (1 - \beta^2)} \left[ (\varepsilon_r \mu_r - \beta^2) \vb{H} - (\varepsilon_r \mu_r - 1) \left( \vb{\beta} \times c \, \vb{D} + \vb{\beta} \left( \vb{\beta} \cdot \vb{H} \right) \right) \right].
	\end{align*}
\end{example}

\subsection{Wigner转动与Thomas进动}

\subsubsection{Wigner转动的通式}

在相对地面系以速度 \(\vb{\beta}_1 c\) 运动的系中，某物体的运动速度为 \(\vb{\beta}_2 c\)，则记该物体相对地面系的速度为 \(\vb{\beta}_2 \oplus \vb{\beta}_1 \cdot c\)，\(\vb{\beta}_1 \oplus \vb{\beta}_2 \cdot c\) 的定义同理.由相对论速度变换知：
\begin{equation*}
	\vb{\beta}_2 \oplus \vb{\beta}_1 = \dfrac{ \left(1 + \dfrac{\gamma_1}{\gamma_1 + 1} \vb{\beta}_1 \cdot \vb{\beta}_2 \right) \vb{\beta}_1 + \dfrac{1}{\gamma_1} \vb{\beta}_2 }{ 1 + \vb{\beta}_1 \cdot \vb{\beta}_2 }, \quad
	\vb{\beta}_1 \oplus \vb{\beta}_2 = \dfrac{ \left(1 + \dfrac{\gamma_2}{\gamma_2 + 1} \vb{\beta}_1 \cdot \vb{\beta}_2 \right) \vb{\beta}_2 + \dfrac{1}{\gamma_2} \vb{\beta}_1 }{ 1 + \vb{\beta}_1 \cdot \vb{\beta}_2 }.
\end{equation*}

\begin{equation*}
	\gamma = \left(1 - |\vb{\beta}_2 \oplus \vb{\beta}_1|^2 \right)^{-1/2} = \left(1 - |\vb{\beta}_1 \oplus \vb{\beta}_2|^2 \right)^{-1/2} = \gamma_1 \gamma_2 \left(1 + \vb{\beta}_1 \cdot \vb{\beta}_2 \right).
\end{equation*}

考虑 \(\vb{\beta}_2 \oplus \vb{\beta}_1 \cdot c\) 物体对应的变换矩阵，易得：
\begin{equation*}
	&M(\vb{\beta}_2) M(\vb{\beta}_1) =
	\begin{pmatrix}
		\gamma                                        & -\gamma \, (\vb{\beta}_2 \oplus \vb{\beta}_1)^\mathrm{T} \\
		-\gamma \, (\vb{\beta}_1 \oplus \vb{\beta}_2) & \vb{A}_{2 \oplus 1}
	\end{pmatrix}.\\
	&\vb{A}_{2 \oplus 1} = \vb{I} + \gamma_1 \gamma_2 \left(1 + \dfrac{\gamma_1 \gamma_2 (\vb{\beta}_1 \cdot \vb{\beta}_2)}{(\gamma_1 + 1)(\gamma_2 + 1)} \right) \vb{\beta}_2 \vb{\beta}_1 + \dfrac{\gamma_1^2}{\gamma_1 + 1} \vb{\beta}_1 \vb{\beta}_1 + \dfrac{\gamma_2^2}{\gamma_2 + 1} \vb{\beta}_2 \vb{\beta}_2.
\end{equation*}

对结果取逆
\begin{equation*}
	\left[ M(\vb{\beta}_2) M(\vb{\beta}_1) \right]^{-1} = M(-\vb{\beta}_1) M(-\vb{\beta}_2) =
	\begin{pmatrix}
		\gamma                                       & \gamma \, (\vb{\beta}_1 \oplus \vb{\beta}_2)^\mathrm{T} \\
		\gamma \, (\vb{\beta}_2 \oplus \vb{\beta}_1) & \vb{A}_{2 \oplus 1}^\mathrm{T}
	\end{pmatrix}.
\end{equation*}
代入
\small
\begin{equation*}
	\bm{\mathcal{I}} &= M(\vb{\beta}_2) M(\vb{\beta}_1) \left[ M(\vb{\beta}_2) M(\vb{\beta}_1) \right]^{-1} =
	\begin{pmatrix}
		\gamma^2 (1 - |\vb{\beta}_2 \oplus \vb{\beta}_1|^2)                                                                   & \gamma^2 (\vb{\beta}_1 \oplus \vb{\beta}_2)^\mathrm{T} - \gamma (\vb{\beta}_2 \oplus \vb{\beta}_1)^\mathrm{T} \cdot \vb{A}_{2 \oplus 1}^\mathrm{T}   \\
		-\gamma^2 (\vb{\beta}_1 \oplus \vb{\beta}_2) + \gamma \, \vb{A}_{2 \oplus 1} \cdot (\vb{\beta}_2 \oplus \vb{\beta}_1) & -\gamma^2 (\vb{\beta}_1 \oplus \vb{\beta}_2)(\vb{\beta}_1 \oplus \vb{\beta}_2)^\mathrm{T} + \vb{A}_{2 \oplus 1} \cdot \vb{A}_{2 \oplus 1}^\mathrm{T}
	\end{pmatrix}.\\
	\bm{\mathcal{I}} &= \left[ M(\vb{\beta}_2) M(\vb{\beta}_1) \right]^{-1} M(\vb{\beta}_2) M(\vb{\beta}_1) =
	\begin{pmatrix}
		\gamma^2 (1 - |\vb{\beta}_1 \oplus \vb{\beta}_2|^2)                                                                             & -\gamma^2 (\vb{\beta}_2 \oplus \vb{\beta}_1)^\mathrm{T} + \gamma (\vb{\beta}_1 \oplus \vb{\beta}_2)^\mathrm{T} \cdot \vb{A}_{2 \oplus 1}             \\
		\gamma^2 (\vb{\beta}_2 \oplus \vb{\beta}_1) - \gamma \, \vb{A}_{2 \oplus 1}^\mathrm{T} \cdot (\vb{\beta}_1 \oplus \vb{\beta}_2) & -\gamma^2 (\vb{\beta}_2 \oplus \vb{\beta}_1)(\vb{\beta}_2 \oplus \vb{\beta}_1)^\mathrm{T} + \vb{A}_{2 \oplus 1}^\mathrm{T} \cdot \vb{A}_{2 \oplus 1}
	\end{pmatrix}.
\end{equation*}
\normalsize
给出关系式
\begin{align*}
	 & \gamma^2 (1 - |\vb{\beta}_2 \oplus \vb{\beta}_1|^2) = \gamma^2 (1 - |\vb{\beta}_1 \oplus \vb{\beta}_2|^2) = 1                                                                                                                                                                                                        \\
	 & -\gamma^2 (\vb{\beta}_1 \oplus \vb{\beta}_2)(\vb{\beta}_1 \oplus \vb{\beta}_2)^\mathrm{T} + \vb{A}_{2 \oplus 1} \cdot \vb{A}_{2 \oplus 1}^\mathrm{T} = -\gamma^2 (\vb{\beta}_2 \oplus \vb{\beta}_1)(\vb{\beta}_2 \oplus \vb{\beta}_1)^\mathrm{T} + \vb{A}_{2 \oplus 1}^\mathrm{T} \cdot \vb{A}_{2 \oplus 1} = \vb{I} \\
	 & \gamma^2 (\vb{\beta}_1 \oplus \vb{\beta}_2)^\mathrm{T} - \gamma (\vb{\beta}_2 \oplus \vb{\beta}_1)^\mathrm{T} \cdot \vb{A}_{2 \oplus 1}^\mathrm{T} = -\gamma^2 (\vb{\beta}_1 \oplus \vb{\beta}_2) + \gamma \, \vb{A}_{2 \oplus 1} \cdot (\vb{\beta}_2 \oplus \vb{\beta}_1) = 0                                       \\
	 & -\gamma^2 (\vb{\beta}_2 \oplus \vb{\beta}_1)^\mathrm{T} + \gamma (\vb{\beta}_1 \oplus \vb{\beta}_2)^\mathrm{T} \cdot \vb{A}_{2 \oplus 1} = \gamma^2 (\vb{\beta}_2 \oplus \vb{\beta}_1) - \gamma \, \vb{A}_{2 \oplus 1}^\mathrm{T} \cdot (\vb{\beta}_1 \oplus \vb{\beta}_2) = 0
\end{align*}

运用以上等式，计算得
\begin{align*}
	R(\vb{\beta}_2 \oplus \vb{\beta}_1) & = M(\vb{\beta}_2) M(\vb{\beta}_1) M^{-1}(\vb{\beta}_2 \oplus \vb{\beta}_1) \\&=
	\begin{pmatrix}
		\gamma                                        & -\gamma \, (\vb{\beta}_2 \oplus \vb{\beta}_1)^\mathrm{T} \\
		-\gamma \, (\vb{\beta}_1 \oplus \vb{\beta}_2) & \vb{A}_{2 \oplus 1}
	\end{pmatrix}
	\begin{pmatrix}
		\gamma                                       & \gamma \, (\vb{\beta}_2 \oplus \vb{\beta}_1)^\mathrm{T}                                                               \\
		\gamma \, (\vb{\beta}_2 \oplus \vb{\beta}_1) & \vb{I} + \dfrac{\gamma^2}{\gamma + 1} (\vb{\beta}_2 \oplus \vb{\beta}_1)(\vb{\beta}_2 \oplus \vb{\beta}_1)^\mathrm{T}
	\end{pmatrix}
	\\&=
	\begin{pmatrix}
		1 & 0                                                                                                                                  \\
		0 & \vb{A}_{2 \oplus 1} - \dfrac{\gamma^2}{\gamma + 1} (\vb{\beta}_1 \oplus \vb{\beta}_2)(\vb{\beta}_2 \oplus \vb{\beta}_1)^\mathrm{T}
	\end{pmatrix} =
	\begin{pmatrix}
		1 & 0                                                     \\
		0 & e^{\theta \, \uv{n}_{2 \oplus 1} \cdot \vb{\epsilon}}
	\end{pmatrix} = e^{-\theta \, \uv{n}_{2 \oplus 1} \cdot \vb{S}}.
\end{align*}
其中
\begin{equation*}
	\theta = \arccos\left( \dfrac{(\gamma + \gamma_1 + \gamma_2 + 1)^2}{(\gamma_1 + 1)(\gamma_2 + 1)(\gamma + 1)} - 1 \right) = 2 \arctan\left( \dfrac{\gamma_1 \gamma_2 |\vb{\beta}_1 \times \vb{\beta}_2|}{1 + \gamma_1 + \gamma_2 + \gamma} \right),\quad
	\uv{n}_{2 \oplus 1} = \dfrac{\vb{\beta}_2 \times \vb{\beta}_1}{|\vb{\beta}_2 \times \vb{\beta}_1|}.
\end{equation*}

另外，也可以证明：
\begin{equation*}
	R(\vb{\beta}_2 \oplus \vb{\beta}_1) = M^{-1}(\vb{\beta}_1 \oplus \vb{\beta}_2) M(\vb{\beta}_2) M(\vb{\beta}_1).
\end{equation*}

这说明，先按 \(\vb{\beta}_1\) 换系、再按 \(\vb{\beta}_2\) 换系，等价于先按 \(\vb{\beta}_2 \oplus \vb{\beta}_1\) 换系、再旋转 \(\theta \, \uv{n}_{2 \oplus 1}\)，也等价于先旋转 \(\theta \, \uv{n}_{2 \oplus 1}\)、再按 \(\vb{\beta}_1 \oplus \vb{\beta}_2\) 换系.

特别地，若 \(\vb{\beta}_1 \perp \vb{\beta}_2\)，记 \(\phi_1 = \artanh\beta_1\)，\(\phi_2 = \artanh\beta_2\)，得：
\begin{equation*}
	\tan\theta = \tanh\left( \dfrac{\phi_1}{2} \right) \tanh\left( \dfrac{\phi_2}{2} \right).
\end{equation*}

特别地，若 \(\vb{\beta}_1 \parallel \vb{\beta}_2\)，则 \(\theta = 0\)，即转角为0.

特别地，若 \(\vb{\beta}_1 = -\vb{\beta}\)，\(\vb{\beta}_2 = \vb{\beta} + \delta\vb{\beta}\)，\(|\delta\vb{\beta}| \ll |\vb{\beta}|\)，则：
\begin{equation*}
	\delta\vb{\phi} = \theta \, \uv{n}_{2 \oplus 1} = \dfrac{\gamma^2}{\gamma + 1} \vb{\beta} \times \delta\vb{\beta}, \quad
	\Delta\vb{\beta} = \vb{\beta}_2 \oplus \vb{\beta}_1 = \gamma \delta\vb{\beta} + \dfrac{\gamma^3}{\gamma + 1} (\delta\vb{\beta} \cdot \vb{\beta}) \vb{\beta},\quad
	M(\vb{\beta} + \delta\vb{\beta}) M(-\vb{\beta}) = R(\delta\vb{\phi}) M(\Delta\vb{\beta}).
\end{equation*}
无穷小速度变缓中的转动亦称为托马斯进动.

\subsubsection{无穷小Wigner转动（托马斯进动）的另一种推导}

\begin{notes}
	考虑 \(f(\lambda) = e^{\lambda(L + \delta L)} e^{-\lambda L}\)，其中 \(\lambda\) 是常量，\(L\) 是算符，\(\delta L\) 是无穷小算符.

	由 \(\d e^{\lambda L}/\d \lambda = L e^{\lambda L} = e^{\lambda L} L\) 得：
	\begin{equation*}
		\dfrac{\d^k f(\lambda)}{\d \lambda^k} = e^{\lambda(L + \delta L)} \left[ L^{k-1}, \delta L \right] e^{-\lambda L}.
	\end{equation*}
	其中定义 \(\left[ L^k, \delta L \right] = \left[ L, \left[ L^{k-1}, \delta L \right] \right]\)，\(\left[ L^0, \delta L \right] = \delta L\)，\(k \in \mathbb{Z}^+\)，\([\cdot\,, \cdot]\) 是对易符号.进而：
	\begin{equation*}
		f(\lambda) = \bm{\mathcal{I}} + \sum_{k=1}^\infty \dfrac{\lambda^k}{k!} \left. \dfrac{\d^k f(\lambda)}{\d \lambda^k} \right|_{\lambda=0} = \bm{\mathcal{I}} + \sum_{k=1}^\infty \dfrac{\lambda^k}{k!} \left[ L^{k-1}, \delta L \right].
	\end{equation*}
\end{notes}

考虑速度 \(\vb{\beta} + \delta\vb{\beta}\)（\(|\delta\vb{\beta}| \ll |\vb{\beta}|\)）的洛伦兹变换：

记 \(\vb{\beta}\) 对应 \(\vb{\phi}\)，\(\vb{\beta} + \delta\vb{\beta}\) 对应 \(\vb{\phi} + \delta\vb{\phi}\)，则：
\begin{equation*}
	\vb{\phi} &= \vb{\beta} \cdot \dfrac{\artanh\beta}{\beta},\\
	\delta\vb{\phi} &= (\vb{\beta} + \delta\vb{\beta}) \cdot \dfrac{\artanh|\vb{\beta} + \delta\vb{\beta}|}{|\vb{\beta} + \delta\vb{\beta}|} - \vb{\beta} \cdot \dfrac{\artanh\beta}{\beta} = \dfrac{\artanh\beta}{\beta} \left( \delta\vb{\beta} - \dfrac{\vb{\beta}(\vb{\beta} \cdot \delta\vb{\beta})}{\beta^2} \right) + \dfrac{1}{1 - \beta^2} \cdot \dfrac{\vb{\beta}(\vb{\beta} \cdot \delta\vb{\beta})}{\beta^2}.
\end{equation*}

记 \(L = \vb{\phi} \cdot \vb{K}\)，\(\delta L = \delta\vb{\phi} \cdot \vb{K}\)，则易证，对于 \(n \in \mathbb{Z}^+\)，有：
\begin{equation*}
	\left[ L^{2n-1}, \delta L \right] &= \dfrac{1}{\beta^2} (\artanh\beta)^{2n} (\vb{\beta} \times \delta\vb{\beta}) \cdot \vb{S}.\\
	\left[ L^{2n}, \delta L \right] &= \dfrac{1}{\beta} (\artanh\beta)^{2n+1} \left( \delta\vb{\beta} - \dfrac{\vb{\beta}(\vb{\beta} \cdot \delta\vb{\beta})}{\beta^2} \right).
\end{equation*}

运用引理的结论，经过一定运算得：
\begin{equation*}
	B\left( \vb{\phi} + \delta\vb{\phi} \right) B^{-1}\left( \vb{\phi} \right)
	= e^{-(\vb{\phi} + \delta\vb{\phi}) \cdot \vb{K}} \, e^{\vb{\phi} \cdot \vb{K}}
	= e^{-(L + \delta L)} \, e^L
	= \bm{\mathcal{I}} - \Delta\vb{\beta} \cdot \vb{K} - \delta\vb{\phi} \cdot \vb{S}
	= R\left( \delta\vb{\phi} \right) B\left( \Delta\vb{\phi} \right).
\end{equation*}
其中
\begin{equation*}
	\delta\vb{\phi} = \dfrac{\gamma^2}{\gamma + 1} \, \vb{\beta} \times \delta\vb{\beta}, \quad
	R\left( \delta\vb{\phi} \right) = e^{-\delta\vb{\phi} \cdot \vb{S}} = \bm{\mathcal{I}} - \delta\vb{\phi} \cdot \vb{S},\quad
	\Delta\vb{\beta} = \gamma \delta\vb{\beta} + \dfrac{\gamma^3}{\gamma + 1} \, (\delta\vb{\beta} \cdot \vb{\beta}) \vb{\beta},
\end{equation*}
\(\Delta\vb{\phi}\) 是 \(\Delta\vb{\beta}\) 的对应量，两者均为小量，故有 \(\Delta\vb{\phi} = \Delta\vb{\beta}\)，且
\begin{equation*}
	B\left( \Delta\vb{\phi} \right) = \bm{\mathcal{I}} - \Delta\vb{\beta} \cdot \vb{K}.
\end{equation*}

易知，上式亦可表示为：

\begin{equation*}
	M\left( \vb{\beta} + \delta\vb{\beta} \right) \, M\left( -\vb{\beta} \right)
	= R\left( \delta\vb{\phi} \right) \, M\left( \Delta\vb{\beta} \right).
\end{equation*}
不难证明，\(\Delta\vb{\beta}\) 恰是 \(\vb{\beta} + \delta\vb{\beta}\) 系相对 \(\vb{\beta}\) 系的速度，\(\delta\vb{\phi}\) 则代表 \(\vb{\beta} + \delta\vb{\beta}\) 系相对 \(\vb{\beta}\) 系的无穷小Wigner转动.

\subsubsection{BMT方程}
在自旋角动量的演化方程（BMT方程）中，也不难提取出托马斯进动角速度.

引入粒子的自旋角动量四维矢量 \(S^\mu\).在粒子本征系中，\(S^\mu\) 的空间分量为本征系中的自旋角动量 \(\vb{s}\)、时间分量为0，进而在地面系中，\(S^\mu\) 的空间分量为地面系中的自旋角动量 \(\vb{S}\)、时间分量为 \(\vb{\beta} \cdot \vb{S}\)（符合\(U_\mu S^\mu = 0\)）.根据洛伦兹变换矩阵，得到 \(\vb{S}\) 与 \(\vb{s}\) 的关系为：
\begin{equation*}
	\vb{S} = \vb{s} + \dfrac{\gamma^2}{\gamma + 1} \left( \vb{\beta} \cdot \vb{s} \right) \vb{\beta}, \quad
	\vb{s} = \vb{S} - \dfrac{\gamma}{\gamma + 1} \left( \vb{\beta} \cdot \vb{S} \right) \vb{\beta}, \quad
	\vb{\beta} \cdot \vb{S} = \gamma \, \vb{\beta} \cdot \vb{s}.
\end{equation*}

对电荷量为\(q\)、静质量为\(m\)、自旋朗德因子为\(g_s\)、以速度 \(\vb{\beta} c\) 运动的粒子，在其瞬时静止参考系 \(K'\)（注意不是本征系，虽然两者时间元同为 \(\d \tau = \gamma^{-1} \d t\)）中，自旋角动量的进动方程为：

\begin{equation*}
	\vb{\mu}_s = \dfrac{g_s q}{2m} \vb{s}, \quad
	\left( \dfrac{\d \vb{s}}{\d \tau} \right)_{K'} = \vb{\mu}_s \times \vb{B}' = \dfrac{g_s q}{2m} \vb{s} \times \vb{B}',\quad
	\vb{B}' = \gamma \left( \vb{B} - \vb{\beta} \times \dfrac{\vb{E}}{c} \right) - \dfrac{\gamma^2}{\gamma + 1} \vb{\beta} \left( \vb{\beta} \cdot \vb{B} \right).
\end{equation*}

将该式写成四维协变形式，左侧显然为 \(\d S^\mu/\d \tau\)，右侧可采用 \(S^\mu, F^{\mu\nu}, U^\mu,\d U^\mu/\d \tau\) 构造.假设右侧仅含电磁场的线性项，注意到 \(F^{\mu\nu} = -F^{\nu\mu}\)，\(\d U^\mu/\d \tau \sim F^{\mu\nu} U_\nu\)，故只能含有 \(F^{\mu\nu} S_\nu\)、\(S_\lambda F^{\lambda\nu} U_\nu U^\mu\)、\(U^\mu S_\nu \d U^\nu/\d \tau\).考虑到 \(\d(U_\mu S^\mu)/\d \tau = 0\) 及其退化形式，可确定各项系数，最终得到如下形式：

\begin{equation*}
	\dfrac{\d S^\mu}{\d \tau} = \dfrac{g_s q}{2m} \left( F^{\mu\nu} S_\nu + \dfrac{1}{c^2} S_\lambda F^{\lambda\nu} U_\nu U^\mu \right) - \dfrac{1}{c^2} U^\mu S_\nu \dfrac{\d U^\nu}{\d \tau} = T^\mu - \dfrac{1}{c^2} U^\mu S_\nu \dfrac{\d U^\nu}{\d \tau},
\end{equation*}
式中
\begin{equation*}
	T^\mu = \dfrac{g_s q}{2m} \left( F^{\mu\nu} S_\nu + \dfrac{1}{c^2} S_\lambda F^{\lambda\nu} U_\nu U^\mu \right).
\end{equation*}
\begin{notes}
	由于通常只考虑粒子在电磁作用力下的运动，即 \(m \,\d U^\mu/\d \tau = q F^{\mu\nu} U_\nu\)，故以上方程常写作：
	\begin{equation*}
		\dfrac{\d S^\mu}{\d \tau} = \dfrac{q}{m} \left( \dfrac{g_s}{2} F^{\mu\nu} S_\nu + \left( \dfrac{g_s}{2} - 1 \right) \dfrac{1}{c^2} S_\lambda F^{\lambda\nu} U_\nu U^\mu \right).
	\end{equation*}
\end{notes}

代入速度四维矢量和自旋角动量四维矢量的具体形式得：
\begin{equation*}
	S_{\nu}\dfrac{\d U^\nu}{\d \tau} = \left( \vb{\beta} \cdot \vb{S} \right) \dfrac{\d(\gamma c)}{\d \tau} - \vb{S} \cdot \dfrac{\d(\gamma \vb{\beta} c)}{\d \tau} = -\gamma c \, \vb{S} \cdot \dfrac{\d \vb{\beta}}{\d \tau}.
\end{equation*}
记 \(T^\mu\) 的空间分量为 \(\vb{T}\)、时间分量为 \(T^0\)，则进动方程的各分量式为：
\begin{equation*}
	\dfrac{\d \vb{S}}{\d \tau} = \vb{T} + \gamma^2 \left( \vb{S} \cdot \dfrac{\d \vb{\beta}}{\d \tau} \right) \vb{\beta}, \quad
	\dfrac{\d(\vb{\beta} \cdot \vb{S})}{\d \tau} = T^0 + \gamma^2 \, \vb{S} \cdot \dfrac{\d \vb{\beta}}{\d \tau}.
\end{equation*}
代入 \(\vb{S}\) 与 \(\vb{s}\) 的关系式，得：
\begin{equation*}
	\dfrac{\d \vb{s}}{\d \tau} = \vb{T} - \dfrac{\gamma}{\gamma + 1} \vb{\beta} T^0 - \dfrac{\gamma^2}{\gamma + 1} \left( \vb{\beta} \times \dfrac{\d \vb{\beta}}{\d \tau} \right) \times \vb{s}.
\end{equation*}
稍加计算，不难得到：
\begin{equation*}
	\vb{T} - \dfrac{\gamma}{\gamma + 1} \vb{\beta} T^0 = -\dfrac{g_s q}{2m} \left[ \gamma \left( \vb{B} - \vb{\beta} \times \dfrac{\vb{E}}{c} \right) - \dfrac{\gamma^2}{\gamma + 1} \vb{\beta} \left( \vb{\beta} \cdot \vb{B} \right) \right] \times \vb{s} = -\dfrac{g_s q}{2m} \vb{B}' \times \vb{s}.
\end{equation*}
故：
\begin{equation*}
	\dfrac{\d \vb{s}}{\d \tau} = \left[ -\dfrac{g_s q}{2m} \vb{B}' - \dfrac{\gamma^2}{\gamma + 1} \left( \vb{\beta} \times \dfrac{\d \vb{\beta}}{\d \tau} \right) \right] \times \vb{s} = \left( \dfrac{\d \vb{s}}{\d \tau} \right)_{K'} - \dfrac{\gamma^2}{\gamma + 1} \left( \vb{\beta} \times \dfrac{\d \vb{\beta}}{\d \tau} \right) \times \vb{s}.
\end{equation*}

可见，\(|\vb{s}|\) 不变，这符合本征系中自旋角动量大小恒定的性质，同时，本征系相对瞬时静止参考系存在旋转，当粒子速度变化 \(\delta\vb{\beta}\) 时，转角
\begin{equation*}
	\delta\vb{\phi} = \dfrac{\gamma^2}{\gamma + 1} \vb{\beta} \times \delta\vb{\beta}.
\end{equation*}
恰是上所求得的无穷小Wigner转动.

下考虑粒子在电磁场中运动，不受机械外力，假设磁场不均匀性与粒子磁矩耦合所引起的梯度力可忽略，即粒子运动方程为 \(m \,\d U^\mu/\d \tau = q F^{\mu\nu} U_\nu\)，则
\begin{equation*}
	m \dfrac{\d(\gamma \vb{\beta} c)}{\d \tau} = q\gamma \left( \vb{E} + \vb{\beta} c \times \vb{B} \right)\quad \implies\quad \dfrac{\d \vb{\beta}}{\d \tau} = \dfrac{q}{mc} \left( \vb{E} - \vb{\beta} (\vb{\beta} \cdot \vb{E}) + \vb{\beta} \times \vb{B} c \right).
\end{equation*}
进而
\begin{align*}
	\dfrac{\d \vb{s}}{\d t}                   & = -\dfrac{g_s q}{2m\gamma} \vb{B}' \times \vb{s} - \dfrac{\gamma^2}{\gamma + 1} \left( \vb{\beta} \times \dfrac{\d \vb{\beta}}{\d t} \right) \times \vb{s}                                                                                                                                                                                                                                           \\&= \dfrac{q}{m} \vb{s} \times \left[ \left( \dfrac{g_s}{2} - 1 + \dfrac{1}{\gamma} \right) \vb{B} - \left( \dfrac{g_s}{2} - 1 \right) \dfrac{\gamma}{\gamma + 1} \vb{\beta} (\vb{\beta} \cdot \vb{B}) - \dfrac{1}{c} \left( \dfrac{g_s}{2} - \dfrac{\gamma}{\gamma + 1} \right) \vb{\beta} \times \vb{E} \right].\\
	\dfrac{\d(\uv{\beta} \cdot \vb{s})}{\d t} & = \uv{\beta} \cdot \dfrac{\d \vb{s}}{\d t} + \dfrac{1}{\beta} \left( \vb{s} - \uv{\beta} (\uv{\beta} \cdot \vb{s}) \right) \cdot \dfrac{\d \vb{\beta}}{\d t} = -\dfrac{q}{m} \left( \vb{s} - \uv{\beta} (\uv{\beta} \cdot \vb{s}) \right) \cdot \left[ \left( \dfrac{g_s}{2} - 1 \right) \uv{\beta} \times \vb{B} + \left( \dfrac{g_s}{2}\beta - \dfrac{1}{\beta} \right) \dfrac{\vb{E}}{c} \right].
\end{align*}

记 \(\theta\) 为 \(\uv{\beta}\) 与 \(\vb{s}\) 的夹角，\(\uv{n}\) 是 \(\vb{s}\) 垂直于 \(\vb{\beta}\) 部分的单位矢量，则上式给出：
\begin{equation*}
	\dfrac{\d \theta}{\d t} = \dfrac{q}{m} \uv{n} \cdot \left[ \left( \dfrac{g_s}{2} - 1 \right) \uv{\beta} \times \vb{B} + \left( \dfrac{g_s}{2}\beta - \dfrac{1}{\beta} \right) \dfrac{\vb{E}}{c} \right].
\end{equation*}

\begin{example}
	对于在质子库仑场中低速运动的电子
	\begin{equation*}
		q=-e,\,m=m_e,\,g_s\approx 2,\,\gamma \approx 1,\,\vb{E} = \dfrac{e}{4\pi\varepsilon_0 r^3}\vb{r},\,\vb{B} = 0.
	\end{equation*}
	代入得：
	\begin{equation*}
		\dfrac{\d \vb{s}}{\d t} = \left( \dfrac{-e\vb{B}_\text{eff}}{m_e} \cdot \dfrac{-\vb{\beta} \times \vb{E}}{2c} \right) \times \vb{s} = -\dfrac{e^2 \vb{L}}{8\pi\varepsilon_0 m_e^2 c^2 r^3} \times \vb{s}.
	\end{equation*}
	其中 \(\vb{L} = m_e \vb{r} \times \vb{\beta} c\) 为电子的轨道角动量，等效附加磁场为
	\begin{equation*}
		\vb{B}_\text{eff} = -\dfrac{\vb{\beta} \times \vb{E}}{2c} = \dfrac{e \vb{L}}{8\pi\varepsilon_0 m_e c^2 r^3}.
	\end{equation*}

	低速下，\(\vb{s} \approx \vb{S}\) 均表示电子自旋角动量，故等效磁场及其产生的附加能量为：
	\begin{equation*}
		\Delta U_{\text{SL}} = -\vb{\mu}_s \cdot \vb{B}_\text{eff} = \dfrac{e^2 \vb{S} \cdot \vb{L}}{8\pi\varepsilon_0 m_e^2 c^2 r^3}.
	\end{equation*}
	该附加能量即为氢原子精细结构理论中的自旋-轨道耦合能.
\end{example}
